\documentclass[a4paper]{article}
\usepackage{amsfonts}
\usepackage{amsmath}
\usepackage{amssymb}
\usepackage{graphicx}    

\begin{document}
  
\noindent \large Auteur: Abdul Ahmed \\
\noindent \large Studierichting: Informatica\\
\noindent \large Oefeningenreeks 2E nr. 28.5 \\

\medskip

\normalsize

$ e^{(\ln(ex))^2} = ex^{\ln(\sqrt{x})}$\\

$ \Leftrightarrow e^{(\ln e + \ln(x))^2} = ex^{\tfrac{1}{2}\ln x}$\\

$ \Leftrightarrow e^{(1 + \ln x)^2} = ex^{\tfrac{1}{2}\ln x}$\\

$ \Leftrightarrow e^{(\ln x)^2+2\ln x+ 1} = ex^{\tfrac{1}{2}\ln x}$\\

$ \Leftrightarrow e^{(\ln x)^2+2\ln x +1 - 1} = x^{\tfrac{1}{2}\ln x}$\\

$ \Leftrightarrow e^{(\ln x)^2+2\ln x} = x^{\tfrac{1}{2}\ln x}$\\

$ \Leftrightarrow (\ln x)^2+2\ln x = \ln x{^{\tfrac{1}{2}\ln x}}$\\

$ \Leftrightarrow (\ln x)^2+2\ln x = {\dfrac{1}{2}(\ln x})(\ln x)$\\

$ \Leftrightarrow (\ln x)^2+2\ln x = \dfrac{1}{2}(\ln x)^2$\\

$ \Leftrightarrow \dfrac{1}{2}(\ln x)^2 + 2\ln x = 0$\\

Stel $t = \ln x$\\

$ \Rightarrow \dfrac{1}{2}t^2 + 2t = 0$\\

$ \Rightarrow t(\dfrac{1}{2}t + 2) = 0$\\

De oplosssingen voor $t$ zijn dan: \\

$ t\textsubscript{1} = -4$ \\

$ t\textsubscript{2} = 0$ \\

Als we $t$ terug vervangen door $\ln x$ krijgen we: \\

$ \ln x\textsubscript{1} = -4$ \\

$ \ln x\textsubscript{2} = 0$ \\


Vervolgens zijn de $x$-waarden: \\

$ x\textsubscript{1} = e^{-4}$ \\

$ x\textsubscript{2} = 1 $\\

\begin{thebibliography}{999}
%\bibitem{a} Referentie
\end{thebibliography}
\end{document} 
